\section{Conclusion}
\label{sec:conclusion}

We presented \textsc{FinOPD}, a financial multi-agent framework that enables policy-level self-evolution through on-policy self-distillation conditioned on hindsight risk-adjusted PnL, Shapley-weighted credit assignment across specialist agents, and a non-parametric belief store that complements parametric adaptation, coupled with regime-adaptive signal weighting and edge-gated abstention at the decision layer. On 2025 Dow-30 equities under strict post-cutoff evaluation against an identical-protocol suite of real trained models and multi-trade agent baselines, \textsc{FinOPD} attains the highest portfolio Sharpe ratio, cumulative return, and Calmar ratio of any method, with a hindsight-conditioned teacher characterizing the performance ceiling (Sharpe 2.96) toward which self-evolution drives the policy. We report the trade-offs transparently: defensive baselines reach a lower absolute drawdown by under-deploying capital, and the per-asset advantage is heterogeneous, but the portfolio-level risk-adjusted leadership is consistent and reproducible.
% 中文：我们提出了 FinOPD，一个通过以事后风险调整 PnL 为条件的在线策略自蒸馏、跨专家 agent 的 Shapley 加权信用分配、补充参数适应的非参数 belief 存储，以及决策层的 regime 自适应信号加权和 edge-gated 弃权来实现策略级自进化的金融多智能体框架。在 2025 年 Dow-30 股票上以严格 post-cutoff 评估、对由真实训练模型与多笔交易 agent 基线构成的同协议套件，FinOPD 取得最高组合夏普和跨评测窗口最稳定的风险调整表现，同时相对被动市场基准显著降低回撤。我们如实报告权衡：系统以风险调整收益与资本保全为目标而非原始收益，其优势在低波动无方向 regime 中收窄。
